% Options for packages loaded elsewhere
% Options for packages loaded elsewhere
\PassOptionsToPackage{unicode}{hyperref}
\PassOptionsToPackage{hyphens}{url}
\PassOptionsToPackage{dvipsnames,svgnames,x11names}{xcolor}
%
\documentclass[
  letterpaper,
  DIV=11,
  numbers=noendperiod]{scrartcl}
\usepackage{xcolor}
\usepackage{amsmath,amssymb}
\setcounter{secnumdepth}{-\maxdimen} % remove section numbering
\usepackage{iftex}
\ifPDFTeX
  \usepackage[T1]{fontenc}
  \usepackage[utf8]{inputenc}
  \usepackage{textcomp} % provide euro and other symbols
\else % if luatex or xetex
  \usepackage{unicode-math} % this also loads fontspec
  \defaultfontfeatures{Scale=MatchLowercase}
  \defaultfontfeatures[\rmfamily]{Ligatures=TeX,Scale=1}
\fi
\usepackage{lmodern}
\ifPDFTeX\else
  % xetex/luatex font selection
\fi
% Use upquote if available, for straight quotes in verbatim environments
\IfFileExists{upquote.sty}{\usepackage{upquote}}{}
\IfFileExists{microtype.sty}{% use microtype if available
  \usepackage[]{microtype}
  \UseMicrotypeSet[protrusion]{basicmath} % disable protrusion for tt fonts
}{}
\makeatletter
\@ifundefined{KOMAClassName}{% if non-KOMA class
  \IfFileExists{parskip.sty}{%
    \usepackage{parskip}
  }{% else
    \setlength{\parindent}{0pt}
    \setlength{\parskip}{6pt plus 2pt minus 1pt}}
}{% if KOMA class
  \KOMAoptions{parskip=half}}
\makeatother
% Make \paragraph and \subparagraph free-standing
\makeatletter
\ifx\paragraph\undefined\else
  \let\oldparagraph\paragraph
  \renewcommand{\paragraph}{
    \@ifstar
      \xxxParagraphStar
      \xxxParagraphNoStar
  }
  \newcommand{\xxxParagraphStar}[1]{\oldparagraph*{#1}\mbox{}}
  \newcommand{\xxxParagraphNoStar}[1]{\oldparagraph{#1}\mbox{}}
\fi
\ifx\subparagraph\undefined\else
  \let\oldsubparagraph\subparagraph
  \renewcommand{\subparagraph}{
    \@ifstar
      \xxxSubParagraphStar
      \xxxSubParagraphNoStar
  }
  \newcommand{\xxxSubParagraphStar}[1]{\oldsubparagraph*{#1}\mbox{}}
  \newcommand{\xxxSubParagraphNoStar}[1]{\oldsubparagraph{#1}\mbox{}}
\fi
\makeatother

\usepackage{color}
\usepackage{fancyvrb}
\newcommand{\VerbBar}{|}
\newcommand{\VERB}{\Verb[commandchars=\\\{\}]}
\DefineVerbatimEnvironment{Highlighting}{Verbatim}{commandchars=\\\{\}}
% Add ',fontsize=\small' for more characters per line
\usepackage{framed}
\definecolor{shadecolor}{RGB}{241,243,245}
\newenvironment{Shaded}{\begin{snugshade}}{\end{snugshade}}
\newcommand{\AlertTok}[1]{\textcolor[rgb]{0.68,0.00,0.00}{#1}}
\newcommand{\AnnotationTok}[1]{\textcolor[rgb]{0.37,0.37,0.37}{#1}}
\newcommand{\AttributeTok}[1]{\textcolor[rgb]{0.40,0.45,0.13}{#1}}
\newcommand{\BaseNTok}[1]{\textcolor[rgb]{0.68,0.00,0.00}{#1}}
\newcommand{\BuiltInTok}[1]{\textcolor[rgb]{0.00,0.23,0.31}{#1}}
\newcommand{\CharTok}[1]{\textcolor[rgb]{0.13,0.47,0.30}{#1}}
\newcommand{\CommentTok}[1]{\textcolor[rgb]{0.37,0.37,0.37}{#1}}
\newcommand{\CommentVarTok}[1]{\textcolor[rgb]{0.37,0.37,0.37}{\textit{#1}}}
\newcommand{\ConstantTok}[1]{\textcolor[rgb]{0.56,0.35,0.01}{#1}}
\newcommand{\ControlFlowTok}[1]{\textcolor[rgb]{0.00,0.23,0.31}{\textbf{#1}}}
\newcommand{\DataTypeTok}[1]{\textcolor[rgb]{0.68,0.00,0.00}{#1}}
\newcommand{\DecValTok}[1]{\textcolor[rgb]{0.68,0.00,0.00}{#1}}
\newcommand{\DocumentationTok}[1]{\textcolor[rgb]{0.37,0.37,0.37}{\textit{#1}}}
\newcommand{\ErrorTok}[1]{\textcolor[rgb]{0.68,0.00,0.00}{#1}}
\newcommand{\ExtensionTok}[1]{\textcolor[rgb]{0.00,0.23,0.31}{#1}}
\newcommand{\FloatTok}[1]{\textcolor[rgb]{0.68,0.00,0.00}{#1}}
\newcommand{\FunctionTok}[1]{\textcolor[rgb]{0.28,0.35,0.67}{#1}}
\newcommand{\ImportTok}[1]{\textcolor[rgb]{0.00,0.46,0.62}{#1}}
\newcommand{\InformationTok}[1]{\textcolor[rgb]{0.37,0.37,0.37}{#1}}
\newcommand{\KeywordTok}[1]{\textcolor[rgb]{0.00,0.23,0.31}{\textbf{#1}}}
\newcommand{\NormalTok}[1]{\textcolor[rgb]{0.00,0.23,0.31}{#1}}
\newcommand{\OperatorTok}[1]{\textcolor[rgb]{0.37,0.37,0.37}{#1}}
\newcommand{\OtherTok}[1]{\textcolor[rgb]{0.00,0.23,0.31}{#1}}
\newcommand{\PreprocessorTok}[1]{\textcolor[rgb]{0.68,0.00,0.00}{#1}}
\newcommand{\RegionMarkerTok}[1]{\textcolor[rgb]{0.00,0.23,0.31}{#1}}
\newcommand{\SpecialCharTok}[1]{\textcolor[rgb]{0.37,0.37,0.37}{#1}}
\newcommand{\SpecialStringTok}[1]{\textcolor[rgb]{0.13,0.47,0.30}{#1}}
\newcommand{\StringTok}[1]{\textcolor[rgb]{0.13,0.47,0.30}{#1}}
\newcommand{\VariableTok}[1]{\textcolor[rgb]{0.07,0.07,0.07}{#1}}
\newcommand{\VerbatimStringTok}[1]{\textcolor[rgb]{0.13,0.47,0.30}{#1}}
\newcommand{\WarningTok}[1]{\textcolor[rgb]{0.37,0.37,0.37}{\textit{#1}}}

\usepackage{longtable,booktabs,array}
\usepackage{calc} % for calculating minipage widths
% Correct order of tables after \paragraph or \subparagraph
\usepackage{etoolbox}
\makeatletter
\patchcmd\longtable{\par}{\if@noskipsec\mbox{}\fi\par}{}{}
\makeatother
% Allow footnotes in longtable head/foot
\IfFileExists{footnotehyper.sty}{\usepackage{footnotehyper}}{\usepackage{footnote}}
\makesavenoteenv{longtable}
\usepackage{graphicx}
\makeatletter
\newsavebox\pandoc@box
\newcommand*\pandocbounded[1]{% scales image to fit in text height/width
  \sbox\pandoc@box{#1}%
  \Gscale@div\@tempa{\textheight}{\dimexpr\ht\pandoc@box+\dp\pandoc@box\relax}%
  \Gscale@div\@tempb{\linewidth}{\wd\pandoc@box}%
  \ifdim\@tempb\p@<\@tempa\p@\let\@tempa\@tempb\fi% select the smaller of both
  \ifdim\@tempa\p@<\p@\scalebox{\@tempa}{\usebox\pandoc@box}%
  \else\usebox{\pandoc@box}%
  \fi%
}
% Set default figure placement to htbp
\def\fps@figure{htbp}
\makeatother





\setlength{\emergencystretch}{3em} % prevent overfull lines

\providecommand{\tightlist}{%
  \setlength{\itemsep}{0pt}\setlength{\parskip}{0pt}}



 


\KOMAoption{captions}{tableheading}
\makeatletter
\@ifpackageloaded{caption}{}{\usepackage{caption}}
\AtBeginDocument{%
\ifdefined\contentsname
  \renewcommand*\contentsname{Table of contents}
\else
  \newcommand\contentsname{Table of contents}
\fi
\ifdefined\listfigurename
  \renewcommand*\listfigurename{List of Figures}
\else
  \newcommand\listfigurename{List of Figures}
\fi
\ifdefined\listtablename
  \renewcommand*\listtablename{List of Tables}
\else
  \newcommand\listtablename{List of Tables}
\fi
\ifdefined\figurename
  \renewcommand*\figurename{Figure}
\else
  \newcommand\figurename{Figure}
\fi
\ifdefined\tablename
  \renewcommand*\tablename{Table}
\else
  \newcommand\tablename{Table}
\fi
}
\@ifpackageloaded{float}{}{\usepackage{float}}
\floatstyle{ruled}
\@ifundefined{c@chapter}{\newfloat{codelisting}{h}{lop}}{\newfloat{codelisting}{h}{lop}[chapter]}
\floatname{codelisting}{Listing}
\newcommand*\listoflistings{\listof{codelisting}{List of Listings}}
\makeatother
\makeatletter
\makeatother
\makeatletter
\@ifpackageloaded{caption}{}{\usepackage{caption}}
\@ifpackageloaded{subcaption}{}{\usepackage{subcaption}}
\makeatother
\usepackage{bookmark}
\IfFileExists{xurl.sty}{\usepackage{xurl}}{} % add URL line breaks if available
\urlstyle{same}
\hypersetup{
  pdftitle={Arrays with Fixed Margins in Stan},
  pdfauthor={Gidon Cohen},
  colorlinks=true,
  linkcolor={blue},
  filecolor={Maroon},
  citecolor={Blue},
  urlcolor={Blue},
  pdfcreator={LaTeX via pandoc}}


\title{Arrays with Fixed Margins in Stan}
\author{Gidon Cohen}
\date{}
\begin{document}
\maketitle


In this note I describe some approaches to modeling matricies and arrays
with fixed row and column margins in Stan where cell counts must be
positive. To assist with modelling I relax the assumption that cell
counts must be integers. Integer cell counts summing exactly to margin
constraints are obtained in generated quantities using integer versions
sequential sampling algorithms described below. The approaches works by
making allocations sequentially, e.g.~to each cell in turn, where each
allocation ensures that the realisation of an overall table or array
summing to the fixed margin remains possible.

To illustrate some of the advantages of sequential sampling I also
describe a raking approach to realising tables with exact margins.

\begin{enumerate}
\def\labelenumi{\arabic{enumi}.}
\tightlist
\item
  Exact margin approach. This approach is like a two-way stick breaking
  algorithm. It uses ideas from Sequential Importance Sampling (Chen et.
  al 2005) to build an algorithm based on an an unconstrained importance
  space to produce cell counts that ensure that both margin constraints
  are met. Thus, \((R - 1)(C - 1)\) unconstrained importance weight
  parameters (\(\lambda\)) are transformed to RxC cell values which sum
  precisely to row and column margins. Because the transform is
  sequential, the Jacobian is quite straightforward.\\
  The resulting cell counts can be modeled in various ways including for
  example by further transforming the cell counts to the simplex
  underlying the table. In the examples below I further transform the
  cell values to an RxC - 1 simplex representation of the table
  (\(\text{ilr}\)) in order to provide structured information about the
  type of table to be generated.
\end{enumerate}

\subsection{Details of the exact margin
approach}\label{details-of-the-exact-margin-approach}

Frèchet bounds, describe the maximum and minimum value of each cell in a
two-way table with known row and column sums as:

\[ \mathrm{max}\left(0, c_k + r_j - T \right) \le \pi_{jk} \le \mathrm{min} \left( r_j, c_k \right)\]

Given these bounds, @Chen2005 describe a procedure the sequentially
allocate values to cells in a way that constructs only valid two-way
contingency table with fixed margins. Their approach begins the
allocation from the top-left cell in the table (\(\pi_{11}\)). In this
table, with no cells allocated they consider the limits on cell
\(\pi_{11}\):

\[ \mathrm{max}\left(0, c_1 + r_1 - T )\right) \le \pi_{11} \le \mathrm{min} \left( r_1, c_1 \right)\]

This is the only condition that \(\pi_{11}\) needs to satisfy.
Recursively suppose that we have chosen \(\pi_{i1}= \pi^{\prime}_{i1}\).
for \(1 \le i \le k -1\), then the only restriction on \(\pi_{k1}\) is
\[ \mathrm{max}\left( 0, \left(c_1 - \sum_{i=1}^{j-1} \pi_{i1}^{\prime} \right) - \sum_{i=j+1}^{J} r_i\right) \le \pi_{j1} \le \mathrm{min}\left(r_j, c_1 - \sum_{i=1}^{j-1} \pi_{i1}^{\prime}\right)\]

This describes a process for sequentially calculating cell bounds from
the top of the first down to the bottom of that column. If the first
column of a contingency table is fully allocated (and only those cells
are allocated) then the remaining unallocated cells in columns
\(2 \dots K\) form a contingency table of dimensions \(J, K-1\) where
the row margins are \(R_1 - \pi_{1,1}, \dots, R_J - \pi_{1,j}\). Hence,
we can continue the allocation process from the top cell of each column
once the previous column is fully allocated. This gives lower bound
(\(\alpha_{jk}\)) and upper bound equation (\(\beta_jk\)) equations that
can be sequentially calculated, so that
\(\alpha_{j,k} = \mathrm{max}\left( 0, \left(c_k - \sum_{i=1}^{k-1} \pi_{i1}^{\prime} \right) - \sum_{i=k+1}^{m} r_i\right)\)
gives the lower bound of cell \(\pi_{j,k}\) and
\(\beta_{[j,k]} = \mathrm{min}\left(r_k, c_1 - \sum_{i=1}^{k-1} \pi_{i1}^{\prime}\right)\)
gives the upper bound of cell \(\pi_{j,k}\).

\subsubsection{Incorporating SIS-constraints in a Stan
Model}\label{incorporating-sis-constraints-in-a-stan-model}

Take a matrix with \(J\) rows and \(K\) columns, where the row margins
are given in a vector \(R_{[1 \dots J]}\) and the column margins given
in a vector \(C_{[1 \dots K]}\), but the the value of each of the
individual cells \(\pi_{j,k}\), the internal structure of the matrix, is
unknown. I describe a way of transforming \(L = (J-1)(K-1)\) parameters
(\(\lambda_{\left[J-1, K-1]\right]}\)) into the \(J \cdot K\) cell
values, \(\pi_{[1 \dots J,1 \dots K]}\), that are the internal structure
of that matrix that sum to the defined row and column margins.

The transformation operates sequentially starting from \(\pi_{[1,1]}\).
Each unconstrained importance parameter \(\lambda_{[j,k]}\) gives a
corresponding cell value \(\pi_{[j,k]}\) between the lower bound
(\(\alpha_{[j,k]}\)) and upper bound (\(\beta_{[j,k]}\)) possible for
the cell given the row and column margins and any allocations already
made to other cells.

For rows \(1 \dots J-1\) and columns \(1 \dots K-1\) the unconstrained
parameter \(\lambda_{[j,k]}\) can be transformed to a value
\(\pi_{j,k}\) in a fixed range from the lower bound \(\alpha_{[j,k]}\)
to the upper bound of \(\beta_{[j,k]}\) with a transformation involving
the inverse logit function:
\[\pi_{[j,k]} = \alpha_{[j,k]} + \frac{1}{1+e^{-\lambda_{[j,k]}}}(\beta_{[j,k]} - \alpha_{[j,k]})\]

For cells in column \(K\) of rows \(1 \dots J-1\),
\(C_{[1 \dots J-1, K]}\) the values is the given by the remaining row
total \(R_{[1 \dots J-1]}\) left to be allocated. For cells in row \(J\)
the value is given by the remaining column totals
\(C\_{\left[1 \dots J \right]}\) left to be allocated.

This transformation gives cells values which are real numbers greater
than or equal to 0. This continuous version is needed for the
\texttt{stan} modelling, and can be thought of as generating something
like an expected value of the cell. The algorithm can be easily amended
to give cells as integer values, by rounding the transformed value to
the nearest integer:

\[\pi_{j,k} = \left\lfloor \alpha_{[j,k]} + \frac{1}{1+e^{-\lambda_{[j,k]}}}(\beta_{[j,k]} - \alpha_{[j,k]}) \right\rceil\]

Stan requires function to be differentiable, so I use continuous
function to calculate approximations the lower bound \(\alpha_{[j,k]}\)
and the upper bound \(\beta_{[j,k]}\). We adapt the continuous hinge
functions described by Andrew Gelman on his
\href{https://statmodeling.stat.columbia.edu/2017/05/19/continuous-hinge-function-bayesian-modeling/}{blog}.

For the lower bound I approximate \$ \mathrm{max}\left(0, c\_1 + r\_1 -
T )\right) \$ with the function
\[\alpha_[j,k] = \delta \log \left(1 +  e^{\frac{c_1 + r_1 - T}{\delta}} \right)\]
As \(c_1 + r_1 - T\) fall below zero then this function quickly
approximates to zero, when it is greater than zero the function
approximates \(c_1 + r_1 - T\). \(\delta\) is a parameter which
determines how quickly that happens (smaller values give a sharper
approximation to the max function).

Similarly, I estimate the upper bound
\(\mathrm{min} \left( r_1, c_1 \right)\) with the function
\[ -\delta \log \left(e^{\frac{-r_1}{\delta}} + e^{\frac{-c_1}{\delta}} \right)\].

Again \(\delta\) is a parameter which determines how quickly that
happens (smaller values give a sharper approximation to the min
function).

\subsubsection{The Log Determinant of the Jacobian of the
Transform}\label{the-log-determinant-of-the-jacobian-of-the-transform}

Fortunately, many of the Jacobian components are zero, and the
sequential nature of the transform means that the Jacobian matrix is
upper triangle. The elements on the diagonal simplify to a convenient
function of bound widths and importance parameters. If we re-index the
\(L = (J - 1)(K - 1)\) parameters for the upper bound (\(\beta\)) the
lower bound (\(\alpha\)) and importance weight (\(\lambda\)) using
\(l = j + (k - 1)(J - 1)\), then:

\[\ln (\det J_{f-1}(y)) = \ln \left( \left(\beta_1 -\alpha_1 \right) \cdot \textrm{logit}^{-1}(\lambda_1) \cdot (1 - \textrm{logit}^{-1}(\lambda_1))\right) + 
\dots + 
\ln \left( \left(\beta_L -\alpha_L \right) \cdot \textrm{logit}^{-1}(\lambda_L) \cdot (1 - \textrm{logit}^{-1}(\lambda_L))\right)\]

\subsubsection{Dealing with Reduced Degrees of Freedom Arising from
Zeros in the Row or Column
Margins}\label{dealing-with-reduced-degrees-of-freedom-arising-from-zeros-in-the-row-or-column-margins}

In a case of an RxC contingency table with fixed margins where none of
the margins are equal to zero, then there are \((R - 1)(C - 1)\) degrees
of freedom, which defines the number of parameters needed to fully
allocate the table by the sequential sampling. In the case that some of
the margins are equal to zero, then this reduces the degrees of freedom
because all of the cells in any row of column with a margin sum of zero
are required to be zero. To deal with the reduced number of parameters
required, I make the allocations in a reduced dimension matrix (with
zero rows and columns removed), and then increase to the RxC dimensions
by inserting rows and/or columns of zero counts into the cells which
make up the rows and/or columns that sum to zero.

\subsection{Further Extensions}\label{further-extensions}

@Chen2006 and @Dinwoodie2011 describe sequential importance sampling
methods for cases which are more complex cases, particular on sampling
from multi-way tables. I have not yet tried to implement these in Stan
but imagine it should be possible to do so.

\subsection{Representations of the table
simplex}\label{representations-of-the-table-simplex}

One finding from these models is that even when descriptions of
simplexes have wide equivalence the way in which they are parameterised
can make a big difference to Stan.

\subsection{An example: votes for parties and
candidates}\label{an-example-votes-for-parties-and-candidates}

\begin{Shaded}
\begin{Highlighting}[]
\NormalTok{name\_lookup }\OtherTok{\textless{}{-}} \ControlFlowTok{function}\NormalTok{(x)\{}
  \FunctionTok{case\_when}\NormalTok{(}
\NormalTok{    x }\SpecialCharTok{==} \DecValTok{1} \SpecialCharTok{\textasciitilde{}} \StringTok{"con"}\NormalTok{,}
\NormalTok{    x }\SpecialCharTok{==} \DecValTok{2} \SpecialCharTok{\textasciitilde{}} \StringTok{"lab"}\NormalTok{,}
\NormalTok{    x }\SpecialCharTok{==} \DecValTok{3} \SpecialCharTok{\textasciitilde{}} \StringTok{"lib"}\NormalTok{,}
\NormalTok{    x }\SpecialCharTok{==} \DecValTok{4} \SpecialCharTok{\textasciitilde{}} \StringTok{"other"}\NormalTok{,}
\NormalTok{    x }\SpecialCharTok{==} \DecValTok{5} \SpecialCharTok{\textasciitilde{}} \StringTok{"green"}\NormalTok{,}
\NormalTok{    x }\SpecialCharTok{==} \DecValTok{6} \SpecialCharTok{\textasciitilde{}} \StringTok{"SNP"}\NormalTok{,}
\NormalTok{    x }\SpecialCharTok{==} \DecValTok{7} \SpecialCharTok{\textasciitilde{}} \StringTok{"spoil"}
\NormalTok{  )}
\NormalTok{\}}

\NormalTok{n\_diag }\OtherTok{\textless{}{-}} \FunctionTok{min}\NormalTok{(}\FunctionTok{length}\NormalTok{(}\FunctionTok{unique}\NormalTok{(actual\_cell\_60}\SpecialCharTok{$}\NormalTok{row\_no)), }
              \FunctionTok{length}\NormalTok{(}\FunctionTok{unique}\NormalTok{(actual\_cell\_60}\SpecialCharTok{$}\NormalTok{col\_no)))}
\NormalTok{diag\_coords }\OtherTok{\textless{}{-}} \FunctionTok{cbind}\NormalTok{(}\DecValTok{1}\SpecialCharTok{:}\NormalTok{n\_diag }\SpecialCharTok{+} \DecValTok{1}\NormalTok{, }\DecValTok{1}\SpecialCharTok{:}\NormalTok{n\_diag }\SpecialCharTok{+} \DecValTok{1}\NormalTok{)}
\end{Highlighting}
\end{Shaded}

\begin{Shaded}
\begin{Highlighting}[]
\NormalTok{ht\_blank }\OtherTok{\textless{}{-}}\NormalTok{ actual\_cell\_60 }\SpecialCharTok{|\textgreater{}}
  \FunctionTok{mutate}\NormalTok{(}
    \AttributeTok{row =} \FunctionTok{name\_lookup}\NormalTok{(row\_no),}
    \AttributeTok{col =} \FunctionTok{name\_lookup}\NormalTok{(col\_no)}
\NormalTok{  ) }\SpecialCharTok{|\textgreater{}}
  \FunctionTok{arrange}\NormalTok{(col\_no, row\_no) }\SpecialCharTok{|\textgreater{}}
  \FunctionTok{select}\NormalTok{(actual\_cell\_value, row, col) }\SpecialCharTok{|\textgreater{}} 
  \FunctionTok{pivot\_wider}\NormalTok{(}\AttributeTok{names\_from =}\NormalTok{ col, }\AttributeTok{values\_from =}\NormalTok{ actual\_cell\_value) }\SpecialCharTok{|\textgreater{}}
  \FunctionTok{adorn\_totals}\NormalTok{(}\AttributeTok{where =} \FunctionTok{c}\NormalTok{(}\StringTok{"row"}\NormalTok{, }\StringTok{"col"}\NormalTok{)) }\SpecialCharTok{|\textgreater{}}
  \CommentTok{\# rename(" " = row) |\textgreater{}}
  \FunctionTok{mutate}\NormalTok{(}\FunctionTok{across}\NormalTok{(}\SpecialCharTok{{-}}\FunctionTok{c}\NormalTok{(row, Total), \textbackslash{}(x) }\FunctionTok{ifelse}\NormalTok{(row }\SpecialCharTok{==} \StringTok{"Total"}\NormalTok{, x, }\StringTok{""}\NormalTok{))) }\SpecialCharTok{|\textgreater{}}
  \FunctionTok{as\_huxtable}\NormalTok{() }\SpecialCharTok{|\textgreater{}}
  \FunctionTok{set\_font\_size}\NormalTok{(}\DecValTok{8}\NormalTok{) }\SpecialCharTok{|\textgreater{}}
  \FunctionTok{set\_bold}\NormalTok{(}\DecValTok{1}\NormalTok{, everywhere, }\ConstantTok{TRUE}\NormalTok{) }\SpecialCharTok{|\textgreater{}}
  \FunctionTok{set\_bold}\NormalTok{(everywhere, }\DecValTok{1}\NormalTok{, }\ConstantTok{TRUE}\NormalTok{) }\SpecialCharTok{|\textgreater{}}
  \FunctionTok{set\_bottom\_border}\NormalTok{(}\DecValTok{1}\NormalTok{, everywhere, }\FloatTok{0.4}\NormalTok{) }\SpecialCharTok{|\textgreater{}}
  \FunctionTok{set\_bottom\_border}\NormalTok{(}\FunctionTok{final}\NormalTok{(}\DecValTok{1}\NormalTok{), everywhere, }\FloatTok{0.4}\NormalTok{) }\SpecialCharTok{|\textgreater{}}
  \FunctionTok{set\_bold}\NormalTok{(}\FunctionTok{final}\NormalTok{(}\DecValTok{1}\NormalTok{), everywhere, }\ConstantTok{TRUE}\NormalTok{) }\SpecialCharTok{|\textgreater{}}
  \FunctionTok{set\_bold}\NormalTok{(everywhere, }\FunctionTok{final}\NormalTok{(}\DecValTok{1}\NormalTok{), }\ConstantTok{TRUE}\NormalTok{)}\SpecialCharTok{|\textgreater{}}
  \FunctionTok{set\_all\_padding}\NormalTok{(}\DecValTok{3}\NormalTok{)}

\FunctionTok{background\_color}\NormalTok{(ht\_blank)[diag\_coords] }\OtherTok{\textless{}{-}} \StringTok{"\#FFF2AC"}
\FunctionTok{number\_format}\NormalTok{(ht\_blank) }\OtherTok{\textless{}{-}} \FunctionTok{list}\NormalTok{(}\ControlFlowTok{function}\NormalTok{(x) }\FunctionTok{format}\NormalTok{(x, }\AttributeTok{big.mark =} \StringTok{","}\NormalTok{, }\AttributeTok{scientific =} \ConstantTok{FALSE}\NormalTok{)) }

\NormalTok{ht\_blank}
\end{Highlighting}
\end{Shaded}





\end{document}
